\providecommand{\cA}{\mathcal{A}}
\providecommand{\cO}{\mathcal{O}}
\providecommand{\paramod}[1]{\Gamma_{#1}}
\providecommand{\paramodlvl}[2]{\Gamma_{#1}(#2)}
\providecommand{\kBKM}{\kappa_{\mathrm{BKM}}}
\providecommand{\kcat}{\kappa_{\mathrm{cat}}}
\providecommand{\kch}{\kappa_{\mathrm{ch}}}
\providecommand{\kfiber}{\kappa_{\mathrm{fiber}}}
\providecommand{\Sp}{\mathrm{Sp}}
\providecommand{\Z}{\mathbb{Z}}
\providecommand{\H}{\mathbb{H}}
\providecommand{\bC}{\mathbb{C}}
\providecommand{\Div}{\mathrm{Div}}
\providecommand{\wt}{\mathrm{wt}}

\section{Humbert normal slices for the Gritsenko-Cl\'ery diagonal divisor}
\label{wn:sec:arnold-boundary-GC}

\subsection*{The two automorphic index sets}

Let
\[
  \cA_2=\Sp_4(\Z)\backslash\H_2 .
\]
Gritsenko-Cl\'ery's diagonal-divisor theorem gives eight genus-two
paramodular forms whose divisor is the diagonal Humbert divisor
\[
  H_1=\left\{
    \begin{pmatrix}\tau&0\\0&\omega\end{pmatrix}
    :\tau,\omega\in\H
  \right\}.
\]
The catalogue is indexed by paramodular triples \((t,N;k)\).  Its row
weights, character orders, and Borcherds constants are
\[
\begin{array}{c|c|c|c|c|c}
r&F_r&(t,N)&\wt(F_r)&m_r&c_r^{\mathrm{dd}}(0)\\ \hline
1&\Delta_5&(1,1)&5&2&10\\
2&M_2(\paramod{2},\nu_4)&(2,1)&2&4&4\\
3&M_3(\paramodlvl{1}{2},\nu_2)&(1,2)&3&2&6\\
4&M_1(\paramod{3},\nu_6)&(3,1)&1&6&2\\
5&M_2(\paramodlvl{1}{3},\nu_2)&(1,3)&2&2&4\\
6&M_{1/2}(\paramod{4},\nu_8)&(4,1)&1/2&8&1\\
7&M_{3/2}(\paramodlvl{1}{4},\nu_4)&(1,4)&3/2&4&3\\
8&M_1(\paramodlvl{2}{2},\nu_4)&(2,2)&1&4&2
\end{array}
\]
where \(m_r\) is the order of the multiplier and
\[
  c_r^{\mathrm{dd}}(0)=2\wt(F_r)=(10,4,6,2,4,1,3,2).
\]
The primitive CHL Borcherds ladder uses the distinct order set
\[
  N\in\{1,2,3,4,6\},\qquad
  c_N(0)=(10,8,6,4,2),\qquad
  \kBKM(\Phi_N)=(5,4,3,2,1).
\]
The common row is \(F_1=\Delta_5=\Phi_1\), with
\(\kBKM(\Delta_5)=10/2=5\).  The monic denominator convention is
\[
  D_5=64^{-1}\Delta_5 .
\]

\subsection*{Local chart}

Fix
\[
  Z_0=\operatorname{diag}(i,i)\in H_1\subset\H_2
\]
and write
\[
  Z=Z_0+W,\qquad
  W=
  \begin{pmatrix}w_1&w_{12}\\ w_{12}&w_2\end{pmatrix}.
\]
On the analytic cover \(\H_2\), the selected branch of \(H_1\) is
\(\{w_{12}=0\}\).  The finite stabiliser of \(Z_0\), the multiplier
system of \(F_r\), and the half-integral square-root convention are
quotient data.  The divisor computation is made first on a local
character cover, then descended to the quotient stack.

\begin{theorem}[Normal slice of the diagonal divisor]
\label{wn:thm:arnold-type-boundary-GC}
Let \(F_r\) be one of the eight Gritsenko-Cl\'ery forms.  On a
sufficiently small analytic character cover \(U\) of \(Z_0\), choose the
branch of the preimage of \(H_1\) given by \(x=w_{12}=0\).  Then
\[
  F_r=u_r x,\qquad u_r(0)\neq0 .
\]
Consequently:
\[
\begin{array}{c|c|c|c|c}
r&m_r&\Div(F_r)\text{ on }U&|F_r|^2\text{ normal form}
&F_r^{m_r}\text{ normal-slice class}\\ \hline
1&2&\text{smooth}&\xi^2+\eta^2&A_1\\
2&4&\text{smooth}&\xi^2+\eta^2&A_3\\
3&2&\text{smooth}&\xi^2+\eta^2&A_1\\
4&6&\text{smooth}&\xi^2+\eta^2&A_5\\
5&2&\text{smooth}&\xi^2+\eta^2&A_1\\
6&8&\text{smooth}&\xi^2+\eta^2&A_7\\
7&4&\text{smooth}&\xi^2+\eta^2&A_3\\
8&4&\text{smooth}&\xi^2+\eta^2&A_3
\end{array}
\]
Here \(x=\xi+i\eta\).  The \(A_{m_r-1}\) entry belongs to the
one-variable normal slice of the character-trivialising power
\(F_r^{m_r}\); the form \(F_r\) itself has a smooth divisor germ.
\end{theorem}

\begin{proof}
Gritsenko-Cl\'ery, Theorem~1.4, states that each row has diagonal
divisor \(H_1\) with multiplicity one on the relevant paramodular
quotient and character cover.  After shrinking \(U\), the local ring
\(\cO_{U,Z_0}\) is regular and \(x=w_{12}\) is a reduced local equation
for the selected branch of \(H_1\).  Multiplicity one gives
\[
  F_r\in (x),\qquad F_r\notin (x^2),
\]
so \(F_r/x\) is a unit.  This proves \(F_r=u_rx\).

The divisor germ \(F_r=0\) is therefore smooth.  The squared Petersson
norm satisfies
\[
  |F_r|^2=|u_r|^2|x|^2 .
\]
The factor \(|u_r|^2\) is a positive smooth unit, and the real Morse
lemma gives the normal form \(\xi^2+\eta^2\) in the real normal plane.

The multiplier of \(F_r\) has order \(m_r\), hence \(F_r^{m_r}\) is
single-valued on the corresponding quotient.  On the normal slice
transverse to \(H_1\),
\[
  F_r^{m_r}=u_r^{m_r}x^{m_r}
\]
is right-equivalent to \(x^{m_r}\).  The Milnor algebra is
\[
  \bC\{x\}/(\partial_x x^{m_r})
  =
  \bC\{x\}/(x^{m_r-1}),
\]
of dimension \(m_r-1\).  This is the simple \(A_{m_r-1}\) normal-slice
singularity in the Arnold-Gusein-Zade-Varchenko classification.
\end{proof}

\subsection*{The \(D_4\) requirement}

The \(D_4\) normal form \(x^2y+y^3\) has Milnor number \(4\) and a
three-line tangent cone over \(\bC\).  A proof of \(D_4\) at the
\((2,2)\) row requires an explicit local equation supplying the extra
branches or an equivalent cubic normal form.  The Gritsenko-Cl\'ery
diagonal-divisor theorem supplies one simple \(H_1\)-branch at the
chosen point.  The available theorem therefore gives the smooth
divisor normal form for \(F_8\) and the \(A_3\) class for the
normal-slice power \(F_8^4\).

\subsection*{BKM and \(K3\times E\) constants}

Borcherds' weight theorem gives
\[
  \kBKM(\Psi_f)=\wt(\Psi_f)=c_f(0)/2
\]
for a Borcherds product with fixed input, chamber, character, and
normalisation.  Applied to the CHL ladder this is
\[
  \kBKM(\Phi_N)=(5,4,3,2,1)
  \quad(N=1,2,3,4,6).
\]
Applied row by row to the diagonal-divisor catalogue it reads
\[
  \kBKM(F_r)=\wt(F_r)
  =(5,2,3,1,2,1/2,3/2,1).
\]
These are automorphic weights.  The local Humbert normal form is
governed by the multiplicity-one divisor statement.

For the compact \(K3\times E\) branch,
\[
  \kcat(K3\times E)=0,\qquad
  \kch(K3\times E)=0,\qquad
  \kappa_{\mathrm{ch}}^{\mathrm{Heis}}(K3\times E)=3,
\]
and
\[
  \kBKM(\Delta_5)=5,\qquad
  \kfiber(K3)=24.
\]
Thus the four construction-specific values are
\[
  \{0,3,5,24\}.
\]
An Arnold normal-slice class supplies local vanishing-cycle data.  A
Cartan matrix comparison for the Borcherds-Kac-Moody algebra requires a
map from vanishing cycles at \(Z_0\) to simple roots.

\subsection*{Source anchors}

Gritsenko-Cl\'ery, \emph{The Siegel modular forms of genus \(2\) with
the simplest divisor}, arXiv:0812.3962, Theorems~1.4 and~3.1: the
eight diagonal-divisor forms and the row-wise Borcherds product
weights.  Borcherds, \emph{Automorphic forms with singularities on
Grassmannians}, Invent.\ Math.\ 132 (1998), Theorem~13.3: the weight
formula \(c(0)/2\).  Arnold-Gusein-Zade-Varchenko,
\emph{Singularities of Differentiable Maps}, Vol.~I, Chs.~11-15:
the \(A_\mu\) simple singularities.  The Vol~III Borcherds-weight
theorem and eight-form catalogue theorem fix the CHL constants, the GC
catalogue constants, and the \(K3\times E\) invariant spectrum;
Gritsenko-Nikulin fixes the monic denominator
\(D_5=64^{-1}\Delta_5\).
